\begin{exercice}\label{exo2-1} (\minsyndical)

Le Professeur Gluckenmuhl effectue des mesures d'\'ecoulements dans les fosses nasales
de trois patients diff\'erents. Pour cela, il place sur chacun des patients un d\'ebitm\`etre en 
entr\'ee des narines, qui mesure le d\'ebit d'air $\Phi$, et un capteur de pression dans la gorge 
qui mesure la diff\'erence de pression entre l'entr\'ee et la sortie des fosses $\Delta P$. Il demande
ensuite au patient d'inspirer et rel\`eve la valeur de $\Delta P$ (en Pascals, Pa) pour diff\'erentes valeurs du d\'ebit $\Phi$ (en litres par minute, l/min), espac\'ees r\'eguli\`erement. 
Il remplit ainsi les trois tableaux ci-dessous:\\

\begin{table}[h]
\center
\begin{tabular}{|c|cccccccc|}
\hline
$\Phi$ & 0 & 1 & 2 & 3 & 4 & 5 & 6 & 7 \\ \hline
$\Delta P$ & -15  & 24 & 59  & 99 & 105 & 129 & 161 & 179 \\ \hline
\end{tabular}
\caption{Relev\'e exp\'erimental. Patient 1.}
\end{table}

\begin{table}[h]
\center
\begin{tabular}{|c|cccccccc|}
\hline
$\Phi$ & 0 & 2 & 4 & 6 & 8 & 10 & 12 & 14 \\ \hline
$\Delta P$ & 13 & 2 & 17 & 29 & 56 & 94 & 152 & 197 \\ \hline
\end{tabular}
\caption{Relev\'e exp\'erimental. Patient 2.}
\end{table}

\begin{table}[h]
\center
\begin{tabular}{|c|cccccccc|}
\hline
$\Phi$ & 0 & 1 & 2 & 3 & 4 & 5 & 6 & 7 \\ \hline
$\Delta P$ & -367 & 390 & 181 & -234 & 25 & 167 & -3 & 154 \\ \hline
\end{tabular}
\caption{Relev\'e exp\'erimental. Patient 3.}
\end{table}

Il d\'ecide ensuite de tracer les trois courbes des Fig. \ref{fig:rglin1} et
\ref{fig:rglin2}, o\`u pour chaque patient, on reporte
les valeurs de d\'ebit en abscisses, et de pression en ordonn\'ees.

\begin{figure}[h!] 
   \centering
   \includegraphics[width=8cm]{matlab/patient1.pdf}
%   \includegraphics[width=10cm]{matlab/patient2.pdf}
%  \includegraphics[width=10cm]{matlab/patient3.pdf} 
   \caption{Mesures d'\'ecoulement effectu\'ees par le Pr. Gluckenmuhl. Patient 1.}
   \label{fig:rglin1}
\end{figure}

\begin{figure}[h!] 
   \centering
   %\includegraphics[width=10cm]{matlab/patient1.pdf}
   \includegraphics[width=8cm]{matlab/patient2.pdf}
   \includegraphics[width=8cm]{matlab/patient3.pdf} 
   \caption{Mesures d'\'ecoulement effectu\'ees par le Pr. Gluckenmuhl. Patient 2 (\`a gauche) et 3 (\`a droite).}
   \label{fig:rglin2}
\end{figure}


Le Pr. Gluckenmuhl veut d\'eduire de ces courbes la loi de variation pression / d\'ebit
et les param\`etres associ\'es \`a cette loi pour chaque patient. Il d\'ecide pour cela
de faire de la r\'egression lin\'eaire.

\begin{enumerate}

\item Pour le Patient 1, il cherche $\beta_0$ et $\beta_1$ tels que :

\[
\Delta P_i = \beta_1 \Phi_i + \beta_0 + \varepsilon_i.
\]

Calculer les indicateurs statistiques $\bar{x}_n$, $\bar{y}_n$, $s^2_x$, $s^2_y$ et $c_{xy}$.
Calculer les coefficients $\beta_0$ et $\beta_1$ en utilisant la m\'ethode des moindres carr\'es,
et tracer sur la Fig. \ref{fig:rglin1} la droite associ\'ee.
D\'eterminer la valeur du coefficient de corr\'elation $r_{xy}$ et de l'erreur quadratique moyenne
$\delta^2$. Qu'en d\'eduisez-vous sur la pertinence de la r\'egression dans ce cas ?

\item Pour le Patient 2, il propose aussi de chercher $\beta_0$ et $\beta_1$ tels que :

\[
\Delta P_i = \beta_1 \Phi_i + \beta_0 + \varepsilon_i.
\]

Ce mod\`ele de r\'egression lin\'eaire vous para\^it-t'il appropri\'e ? Justifiez. Que sugg\'erez-vous
\`a la place ? En suivant la m\^eme d\'emarche qu'en 1., d\'eterminer la valeur optimale
des coefficients au sens des moindres carr\'es et tracer la courbe associ\'ee. 
D\'eterminer la valeur du coefficient de corr\'elation $r_{xy}$ et de l'erreur quadratique moyenne
$\delta^2$ et comparer \`a ce qu'on aurait obtenu \`a l'aide du mod\`ele 
$\Delta P_i = \beta_1 \Phi_i + \beta_0 + \varepsilon_i$.

\item Pour le Patient 3, le Pr. Gluckenmuhl utilise aussi le mod\`ele :

\[
\Delta P_i = \beta_1 \Phi_i + \beta_0 + \varepsilon_i.
\]

Il d\'etermine la droite aux moindres carr\'es et d\'ecide, dans un exc\`es d'enthousiasme, 
d'envoyer ses r\'esultats \`a l'Acad\'emie des Sciences. Calculer le coefficient de corr\'elation 
$r_{xy}$ pour le Patient 3, et en d\'eduire si, suite \`a ses travaux , le Pr. Gluckenmuhl recevra le Prix
Nobel ou le Prix Ig-Nobel ({\sf www.ignobel.com}).

\end{enumerate}

{\it Remarque d'ordre culturel:} De telles mesures de relation pression-d\'ebit sont utilis\'ees dans la
pratique pour mieux comprendre certaines pathologies respiratoires, et \'etudier l'impact de dispositifs
m\'edicaux (respirateurs artificiels...). Si vous avez quelques souvenirs de m\'ecanique des fluides, 
remarquez que la relation lin\'eaire pour le Patient 1 peut \^etre caus\'ee par un \'ecoulement \`a faible 
vitesse r\'egi par la loi de Poiseuille (le Patient 1 a certainement le nez bouch\'e). Pour le Patient 2,
on peut supposer un \'ecoulement r\'egi par la formule de Bernoulli (pourquoi ?).\\

%\corrref{TD1_2}

%rg = 

%        phi1: [0 1 2 3 4 5 6 7]
%        phi2: [0 2 4 6 8 10 12 14]
%     deltap1: [0 25 50 75 100 125 150 175]
%     deltap2: [0 4 16 36 64 100 144 196]
%     deltap3: [0 25 50 75 100 125 150 175]
%    deltap1n: [-15 24 59 99 105 129 161 179]
%    deltap2n: [13 2 17 29 56 94 152 197]
%    deltap3n: [-367 390 181 -234 25 167 -3 154]

%LS1 = 

%    phimoy: 3.5000
%     dpmoy: 92.6250
%    phivar: 5.2500
%     dpvar: 3.8845e+03
%    covpdp: 141.1875
%      rpdp: 0.9887
%     beta1: 25
%     beta0: 0
%    beta1e: 26.8929
%    beta0e: -1.5000
%    error2: 87.5491
%    error1: 9.3568
%     phiLS: [1x71 double]
%      dpLS: [1x71 double]

%LS2 = 

%    phimoy: 7
%     dpmoy: 70
%    phivar: 21
%     dpvar: 4.4985e+03
%    covpdp: 287
%      rpdp: 0.9338
%     beta1: 14
%     beta0: 0
%    beta1e: 13.6667
%    beta0e: -25.6667
%    error1: 576.1667
%    error2: 24.0035
%     phiLS: [1x71 double]
%      dpLS: [1x71 double]

%
%LS3 = 

%    phimoy: 70
%     dpmoy: 70
%    phivar: 4452
%     dpvar: 4.4985e+03
%    covpdp: 4451
%      rpdp: 0.9946
%     beta1: 1
%     beta0: 0
%    beta1e: 0.9998
%    beta0e: 0.0157
%    error1: 48.4998
%    error2: 6.9642
%     phiLS: [1x71 double]
%      dpLS: [1x71 double]

%LS4 = 

%    phimoy: 3.5000
%     dpmoy: 39.1250
%    phivar: 5.2500
%     dpvar: 5.1787e+04
%    covpdp: 118.6875
%      rpdp: 0.2276
%     beta1: 25
%     beta0: 0
%    beta1e: 22.6071
%    beta0e: -40
%    error2: 4.9104e+04
%    error1: 221.5946
%     phiLS: [1x71 double]
%      dpLS: [1x71 double]

\end{exercice}
